\section{基本原则}
\begin{enumerate}
  \item 12 项基本原则。刑诉的基本原则 \bdy 刑诉法的基本原则。
  \item 体现刑诉的理念和目的，贯穿刑诉活动的全过程 \hz 主要阶段，对刑诉活动具有指导意义和规范作用，是刑诉的主体都要遵循的基本准则。
  \item 特征。
        \begin{enumerate}
          \item 体现刑事诉讼活动的基本规律。\zd{控权保民}。
          \item 必须由\zd{法律明确规定}（法定性）。一般原则、独有原则。
          \item 贯穿全过程 \hz 主要诉讼阶段。
          \item 具有\zd{法律约束力}。
        \end{enumerate}
  \item \zd{侦查权、检察权、审判权}由专门机关依法行使。
        \begin{enumerate}
          \item 权力分工。侦查权主要分给公安（国安、军队、监狱、海警局、检）。检察权检察院。审判权法院。
          \item 权力专属。
          \item 依法。行使各自的权力都要依据刑法、刑诉法来行使。
        \end{enumerate}
  \item 严格遵守法律程序/\zd{程序法定}原则。不同阶段违法的后果不同。\al{以事实为依据，以法律为准绳}。
        \begin{enumerate}
          \item 侦查阶段非法搜集证据 \yjt 非法证据排除。
          \item 一审 \yjt 二审应当裁定发回重审。
          \item 死刑复核 \yjt 应当裁定不予核准，发回重审。
          \item 生效裁判 \yjt 应当启动再审进行纠正。
        \end{enumerate}
  \item \san \zd{法院、检察院依法独立行使职权}。\zd{\bdy 三权分立 \bdy 司法独立}。
        \begin{enumerate}
          \item 法院上下级间彼此\zd{独立}。\zd{监督与被监督}。\al{指导、命令、请示、汇报 \cuo}。
          \item 检察院上下级是领导与被领导关系。\al{指导、命令、请示、汇报 \dui}。\zd{个体不独立，整个系统独立}。
          \item 法院和检察院\zd{门当户对}。
        \end{enumerate}
  \item \zx{分工}负责，互相配合，互相制约。\zx{主体仅限于公检法三家}。
  \item \san \zd{检察院依法对刑事诉讼实行法律监督}。\zd{不同阶段监督方式不同}。
        \begin{enumerate}
          \item 监督方式。
                \begin{enumerate}
                  \item 情节较轻 \yjt \zd{口头}提出\zd{纠正意见}。
                  \item 情节较重 \yjt \zd{纠正违法通知书}。\zd{纠正针对的是确有错误的，找同级法院}。
                  \item 违法情形 \yjt 提出\zd{检察建议}。\zd{建议是我不服，找上级法院}。
                  \item 犯罪的 \yjt \zd{追究刑事责任}。
                \end{enumerate}
          \item 立案监督。\yi 当立不立。先要求说理，后通知立案。\er 不当立却立。先要求说理，后通知撤案。
          \item 侦查监督。\bs{立案之后}！！\al{在立案之前商情检察院 \cuo}。
                \begin{enumerate}
                  \item \yi 介入侦查。公安商请 \hz 确有必要。\er 介入调查。监察机关\zd{商请}。
                  \item 提起批捕。\yi 公安漏报。先求报，后自决。\er 自侦部门漏报。先建议，后决定。
                  \item 漏罪、漏人 \yjt 补充侦查 \hz 补充移送 \hz 直接提起公诉。
                \end{enumerate}
          \item 审判监督。
                \begin{enumerate}
                  \item \zd{庭审后}对庭审活动提出纠正意见。
                  \item 对刑事判决、裁定进行\zd{抗诉}。
                        \begin{enumerate}
                          \item \zd{二审}抗诉。一审法院同级检，针对一审\zd{未生效}裁判，向上一级法院提出二审抗诉。
                          \item \zd{再审}抗诉。上级检，针对下级法院\zd{已生效判决}，向同级法院提出再审抗诉。
                        \end{enumerate}
                  \item \zd{死刑复核}。最高检\zd{可以}提出意见，最高法\zd{应当}结果通报。
                  \item 特别程序。决定 \yjt 纠正意见。裁定、判决 \yjt 抗诉。
                \end{enumerate}
          \item 执行监督。死刑 \yjt 临场。监外执行、减刑/假释 \yjt 事前/事后。
        \end{enumerate}
  \item 各民族公民有权使用本民族语言文字进行诉讼。
        \begin{enumerate}
          \item 当事人、其他诉讼参与人。方言 \bdy 本民族语言。有权 = 权利\zd{可以放弃}。
          \item 公检法\zd{应当}使用当地通用语言进行各项活动，\zd{应当聘请翻译}。
        \end{enumerate}
  \item \zd{嫌疑人/被告人}有权获得\zd{辩护。\bdy 辩论}。
        \begin{enumerate}
          \item 告知义务。第一次讯问 \hz 采取强制措施、检自收到移送审查起诉材料 3 日内、法院受案 3 日内。
          \item 保障义务。特殊情况下，没有委托辩护人的，应当通知法援机构指派律师。
        \end{enumerate}
  \item \san \zx{未经法院依法判决，对任何人不得确定有罪}。
        \begin{enumerate}
          \item \zd{\bdy 疑罪从无}。仅一审，二审疑罪可查可发回。\al{可以得出疑罪从无的结论 \cuo}。
          \item \zd{\bdy 无罪推定}。无罪推定是假定无罪，推定无罪。\al{表明已经全面认同和确立无罪推定原则 \cuo}。
          \item 达到定罪标准，方可定罪（排除合理怀疑）。假定无罪的待遇。
          \item 法院定罪必须严格依照法定程序。
          \item \zd{区分称谓，以提起公诉为节点}。
        \end{enumerate}
  \item 保障\zd{诉讼参与人}的诉讼权利。核心是保护嫌疑人/被告人的辩护权。
  \item \yi \zd{认罪认罚从宽}。
        \begin{enumerate}
          \item 认罪。\zx{认全部}。\yi 如实供述，明确表示承认犯罪\zx{事实}。\er 只认一个 \hz 一部分全案都不算。
          \item 认罚。
                \begin{enumerate}
                  \item 接受刑罚的后果。
                  \item 接受程序的简化。\zx{但有选择程序的权利}，可以不接受简易程序、速裁程序。
                  \item 必须要真心实意的认罚\bs{悔}过。不能表里不一、阳奉阴违。\al{串供、毁灭证据、转移财产}。
                  \item \al{甲杀害妻子的情人丙，认罪认罚，但表示并不后悔杀丙，对丙不能适用 \cuo}。
                  \item \al{甲认为被害人也有过错不同意全额赔偿 \dui}。
                  \item \al{甲辩称欺骗乙感情为真，但 8W 款项属于民间借贷并在积极退赔中 \cuo}。
                  \item \zd{结合民法、常识判断辩解合不合理 \yjt 合理 \yjt 构成}。
                \end{enumerate}
          \item 从宽。\yi 量刑优惠。\er 可以更轻的强制措施、宽缓。
          \item \zx{与自首、坦白不作重复评价}。
          \item 价值排序。\zd{早 \dy 晚，主动 \dy 被动，彻底认 \dy 不彻底认，稳定认 \dy 不稳定认}。
          \item 制度价值。
                \begin{enumerate}
                  \item 承载现代司法宽容精神。
                  \item 探索形成非对抗的诉讼格局。减少上诉和申诉的几率，更有利于社会的稳定。
                  \item 公正基础上的效率观，实现司法资源的优化配置。
                  \item \al{提高了嫌疑人、被告人的刑事诉讼地位 \dui{刑诉本来就不是平等两造的诉讼}}。
                  \item \al{体现了职权主义诉讼构造 \cuo{当事人主义}}。
                \end{enumerate}
          \item 适用的范围、阶段\bs{无禁区}。\al{监委调查阶段不能适用认罪认罚制度 \cuo}。
          \item 注意保障被害方权益。
                \begin{enumerate}
                  \item \zd{应当听取}被害人意见。
                  \item 促进和解、谅解。
                  \item \zx{被害人异议不影响适用}。从宽处理并不以被害人同意作为前提。
                \end{enumerate}
          \item \zd{\bdy 免除刑事处罚}。例外。\zd{如实供、大立功、最高检核}。
          \item 不同阶段的适用。\zd{都是应当}！！
                \begin{enumerate}
                  \item 侦查。应当告知，应当记录在案。\al{侦查阶段被逮捕后，应开展羁押必要性审查 \dui}。
                  \item 审查起诉。
                        \begin{enumerate}
                          \item 应当\zd{告知权利与听取意见}（嫌疑人、辩护人 \hz 值班律师、被害人及其诉讼代理人）。
                                \al{没有辩护人的，应当通知值班律师 \dui}。
                          \item 审查真实性、合法性。是否自\zd{愿}、精神正\zd{常}、\zd{理}解后果、\zd{告}知权利、写\zd{明}情况、态度真\zd{诚}，\zd{违遗愿，可重来}。
                          \item 应当签署具结书。例外。\zd{盲聋哑、半疯傻、小孩法辩不认罚}。
                                \begin{itemize}
                                  \item 辩护人 \hz 值班律师在场，未成年法定代理人。
                                  \item 签了具有约束力。
                                \end{itemize}
                          \item 检察院应当\zd{提出量刑建议}。\yi 量刑建议书 \hz 起诉书写明。\er 原则确定型，新类型可幅度。
                                \al{可自行对甲进行社会调查 \dui}。
                        \end{enumerate}
                  \item 审判。
                        \begin{enumerate}
                          \item 审查真实性、合法性。
                          \item \zd{量刑建议的采纳与调整}。\zd{一般应当}。例外时可以不采纳量刑建议。
                                \begin{itemize}
                                  \item 质的问题。\zd{无罪无责违意愿，否认、翻供、罪}名\zd{不一}致或其他。
                                  \item 量的问题。\zd{量刑建议显不当，应当建议检调整，仍然不当依法判} \yjt
                                        检察院如果认为量刑不当，提抗诉、检察建议、纠正违法通知书。
                                  \item 调整量刑建议。调。\zd{不调 \yjt 说明理由、依据}。
                                  \item 时间。庭前 \hz 休庭（排除速裁）、控辩双方不在场 \yjt 重新听取意见、庭审 \yjt 记录在案。
                                  \item \zd{最迟法庭审理结束前}。
                                \end{itemize}
                          \item 普通程序简化审。\al{主要围绕量刑进行 \dui}。
                          \item 当庭认。
                          \item 二审认。\sjt 不用退回，但早优于晚。
                        \end{enumerate}
                \end{enumerate}
          \item 反悔。\zd{一夜回到认罪前，从宽待遇要收回}。
                \begin{enumerate}
                  \item 起诉前。具结书\zd{失效}。
                  \item 酌定不起诉后。法定不起诉 \yjt \zd{撤销重做}。酌定不起诉 \yjt 维持。依法\zd{起诉}。
                  \item 审理中。该咋判咋判。检。\yi \zd{撤回}。\er \zx{建议}法院适用普通程序、重新考虑量刑。
                        \al{法院当庭对甲在量刑方面从重处罚 \cuo}。
                  \item 裁判后。应当依法提出抗诉。
                \end{enumerate}
        \end{enumerate}
  \item \er \zd{具有法定情形不予追究刑事责任}。\zx{显著轻，过时效，特赦告诉和死掉}。\zd{只针对存在违法犯罪事实的情形}。
        \begin{enumerate}
          \item 情节\zd{显著}轻微、危害不大，\zd{不认为是犯罪}。违法了\zd{还不构犯罪，\bdy 犯罪情节轻微}。
                \al{甲买水果与人争执，水果刀刺中乙后背。事后乙被鉴定为轻微伤。检认为甲的行为没有造成严重后果，情节显著轻微，不认为是犯罪 \dui}。
          \item 犯罪\zd{已过追诉时效}。最高刑不满 5 年 \yjt 5 年。5 - 10 \yjt 10。10+ \yjt 15。无期、死刑 \yjt 20（特别追诉最高检批准）。
                特别追诉 = 未同意，不可以提起公诉。\al{经审理发现犯罪已过追诉时效}。
          \item 经特赦令免除刑罚的。
          \item 告诉才处理的。没有 \hz 撤回告诉。
          \item 嫌疑人/被告人\zd{死}亡的。\zd{仅限自然人，单位“死”不算}。
          \item 其他。\al{强疗}。
          \item \zd{不同阶段不同的处理方式}。\zd{忽略题干其他描述}。\al{被害人要求甲赔偿丧葬费等损失，后甲死亡}。
                \begin{enumerate}
                  \item 立案 \yjt 不立案决定。
                  \item 侦查 \yjt 撤案决定。\al{公安侦查中发现甲神情呆滞，经鉴定属于依法不负刑事责任的精神病人}。
                  \item 审查起诉 \yjt 不起诉决定。
                  \item 审判 \yjt 终止审理裁定 \yjt If(无罪) \yjt 判决。\al{据以查明的事实和证据，尚不能确认其无罪}。
                \end{enumerate}
        \end{enumerate}
  \item 追究外国人刑事责任适用我国刑诉法。例外。外交特权和豁免权通过外交途径解决。
\end{enumerate}
